% Generated by generate.py; do not edit.
\begin{figure}
\begin{tikzpicture}[metamodel]
\node[class] (Exam) {\begin{tabular}{l}\multicolumn{1}{c}{\textbf{Exam}}\\\hline name: str\\weight: int\\\hline \textit{inv} exam\_has\_name\end{tabular}};
\node[class] (Student) {\begin{tabular}{l}\multicolumn{1}{c}{\textbf{Student}}\\\hline name: str\\\hline \textit{inv} student\_has\_name\end{tabular}};
\node[class] (TimeSlot) {\begin{tabular}{l}\multicolumn{1}{c}{\textbf{TimeSlot}}\\\hline id: int\\\hline \textit{inv} time\_slot\_has\_non\_negative\_id\end{tabular}};
\path (Exam) edge[association, "exams 1..*" source role, "students 0..*" target role] (Student);
\path (Exam) edge[association, "exams 0..*" source role, "timeSlot 1" target role] (TimeSlot);
\end{tikzpicture}
\caption{The Problem DSL metamodel of the BESSER chain, drawn from \texttt{besser/output/metamodels/problem\_dsl.json}, which the chain exports from \texttt{metamodels/buml/problem\_dsl.py}. Part~I makes \texttt{timeSlot} 1; the chain checks it only on the solved problem model, since an unsolved problem has no slots. T1 does not use \texttt{weight}: the paper has no \term{exam} weights.}\label{fig:metamodel-besser-problem-dsl}
\end{figure}
